% =====================================================================
% Results
% Status: placeholder — no empirical numbers are entered until the
% pipeline of python_project/notebooks/ has been run.
% Each subsection includes a "what will be reported" template so
% that the structure is fixed before the analysis is done.
% =====================================================================
\section{Results}\label{sec:results}

This section reports the empirical estimates produced by the
pipeline of Section~\ref{sec:methodology} on the main sample of
Section~\ref{sec:emp-samples}. Section~\ref{sec:results-marginals}
reports the marginal extremes per station, season by season.
Section~\ref{sec:results-pairwise} reports the pairwise
dependence estimates. Section~\ref{sec:results-distance} reports
the dependence--distance summary, and
Section~\ref{sec:results-comparison} reports the seasonal
comparison.

% TODO: every numerical claim below is to be inserted only after
%       the pipeline has been executed. Do not pre-write
%       conclusions.

\subsection{Marginal extremes per
station}\label{sec:results-marginals}

We first inspect the marginal upper tail of \(\mathrm{FX}\) at
each station, separately for winter and summer.

\paragraph{What is reported.}
\begin{itemize}
    \item Empirical 95\%, 97\% and 99\% within-(station, season)
          quantiles of hourly \(\mathrm{FX}\).
    \item GEV parameter estimates
          \((\widehat\mu_i^{(c)}, \widehat\sigma_i^{(c)},
          \widehat\xi_i^{(c)})\) on the seasonal block-maxima
          series, with the standard sandwich standard errors.
    \item GPD parameter estimates
          \((\widehat\sigma_{u,i}^{(c)}, \widehat\xi_i^{(c)})\) at
          the chosen tail level (Section~\ref{sec:emp-extremes}).
    \item Goodness-of-fit diagnostics: return-level plot and QQ
          plot per (station, season).
\end{itemize}

% TODO: insert Figure \ref{fig:marginal-fits-per-station}
%       (\texttt{marginal\_fits\_per\_station.pdf}) and Table
%       \ref{tab:marginal-fits} (\texttt{marginal\_fits.tex}) once
%       produced by notebook 02.

\paragraph{Interpretation template.}
After inspection, this paragraph will report whether the shape
parameter \(\widehat\xi_i^{(c)}\) is systematically different
between winter and summer, and whether the seasonal contrast in
the marginal upper tail justifies the use of a season-conditional
marginal transform in Section~\ref{sec:results-pairwise}.

\subsection{Pairwise extremal
dependence}\label{sec:results-pairwise}

We now estimate the pairwise dependence summaries
\(\widehat\chi_{ij}^{(c)}(u)\), \(\widehat\theta_{ij}^{(c)}\) and
\(\widehat{\bar\chi}_{ij}^{(c)}(u)\) for every pair on the balanced
panel.

\paragraph{What is reported.}
\begin{itemize}
    \item Scatter plot of
          \(\widehat\chi_{ij}^{(c)}(u_0)\) versus \(d_{ij}\) for
          each season, at the reference level \(u_0\)
          (Section~\ref{sec:emp-extremes}).
    \item Scatter plot of \(\widehat\theta_{ij}^{(c)}\) versus
          \(d_{ij}\) using the F-madogram estimator on block
          maxima.
    \item Joint plot of \(\widehat\chi_{ij}^{(c)}(u)\) and
          \(\widehat{\bar\chi}_{ij}^{(c)}(u)\) on the grid of
          \(u\), for a small set of representative pairs at
          different distances.
\end{itemize}

% TODO: insert Figure
%       \ref{fig:pair-chi-vs-distance-by-season}
%       (\texttt{pair\_chi\_vs\_distance\_by\_season.pdf}) once
%       produced by notebook 03.

\paragraph{Interpretation template.}
The text will note whether the \(\widehat\chi_{ij}^{(c)}(u)\)
estimates appear to converge to a strictly positive limit (the
asymptotic-dependence case) or to decay to zero (the asymptotic-
independence case), and whether the diagnosis differs between
winter and summer.

\subsection{Dependence as a function of
distance}\label{sec:results-distance}

The pair-level estimates are noisy; we summarise them by the
local-linear smoother
\(\widehat\chi^{(c)}(d)\) defined in
Section~\ref{sec:meth-distance}.

\paragraph{What is reported.}
\begin{itemize}
    \item Smoothed curve \(\widehat\chi^{(c)}(d)\) with block-
          bootstrap percentile band, one curve per season, on a
          single panel.
    \item Same plot for \(\widehat\theta^{(c)}(d)\).
    \item Bandwidth used, plus a small bandwidth-sensitivity
          panel.
\end{itemize}

% TODO: insert Figure \ref{fig:chi-distance-curves}
%       (\texttt{chi\_distance\_curves.pdf}) once produced by
%       notebook 03.

\paragraph{Interpretation template.}
The text will describe the qualitative shape of each curve
(monotone decay, intercept near zero distance, behaviour at the
maximum distance in the panel), report the smoother bandwidth, and
note whether the two seasonal curves are visually separable on the
range of distances spanned by the balanced panel.

\subsection{Winter--summer comparison}\label{sec:results-comparison}

Finally we compare the two seasonal curves directly.

\paragraph{What is reported.}
\begin{itemize}
    \item Scalar summaries
          \(\widehat\chi^{(c)}(d_0)\) at the reference distances
          \(d_0 \in \{50, 100, 200\}\) km (or the revised set of
          Section~\ref{sec:emp-distances}), with bootstrap
          percentile intervals, in a two-column (winter, summer)
          table.
    \item Pointwise difference
          \(\widehat\chi^{(\mathrm{W})}(d) -
          \widehat\chi^{(\mathrm{S})}(d)\) with bootstrap band.
    \item Permutation-test \(p\)-value for the null
          \(\chi^{(\mathrm{W})}(d) = \chi^{(\mathrm{S})}(d)\) on
          the reference distance grid.
\end{itemize}

% TODO: insert Figure
%       \ref{fig:chi-distance-winter-minus-summer}
%       (\texttt{chi\_distance\_winter\_minus\_summer.pdf}) and
%       Table \ref{tab:seasonal-summary}
%       (\texttt{seasonal\_summary.tex}) once produced by notebook
%       03.

\paragraph{Interpretation template.}
The text will state, in careful language (``evidence suggests''
rather than ``proves''), the sign and approximate magnitude of the
seasonal contrast at the reference distances, whether the
permutation test rejects the null at the chosen significance
level, and how the contrast varies with distance. Climate-change
interpretations are explicitly avoided. The interpretation will
explicitly distinguish finite-level dependence (i.e.\ what
\(\widehat\chi_{ij}^{(c)}(u)\) reports at the chosen \(u\)) from
the asymptotic-dependence question discussed in
Section~\ref{sec:results-pairwise}.
